\lecture{14}
\title{Icosahedron Symmetries, Jordan-Holder}

\begin{document}

The only group we've shown to be simple is $C_p$ for primes $p$.  Let's find a more complicated simple group.

\subsection{Icosahedron Symmetries are Simple}

Consider $G$ to be the rotational symmetry group of a regular icosahedron.
\begin{center}
    \frame{\includegraphics[width=0.2\textwidth]{lecture-02/regular-icosahedron.png}}
\end{center}
We know from Lecture 12 that $G \subseteq SO_3(\mathbb{R})$ has $60$ group elements.  These elements are:
\[
\text{group elements of } G = 
\begin{cases}
    \text{identity rotations} & (+1) \\
    \text{rotations around an axis through a face} & (+20) \\
    \text{rotations around an axis through an edge} & (+15) \\
    \text{rotations around an axis through a vertex} & (+24)
\end{cases}
\]
What do the conjugacy classes of $G$ look like?  Well, conjugation is just a change of coordinates, so:
\[
\text{conjugacy classes of } G = 
\begin{cases}
    \text{identity rotations} & (+1) \\
    \text{rotations around an axis through a face} & (+20) \\
    \text{rotations around an axis through an edge} & (+15) \\
    \text{rotations by $\pm \frac{2\pi}{5}$ about an axis through a vertex} & (+12) \\
    \text{rotations by $\pm \frac{4\pi}{5}$ about an axis through a vertex} & (+12) \\
\end{cases}
\]
The above will let us show $G$ is simple!

\begin{theorem}{Simplicity of Icosahedra}
    The group $G$ is simple.
\end{theorem}

\begin{proof}
    Suppose $H \unlhd G$ is a normal subgroup.  Recall this means $gHg^{-1} = H$ for all $g \in G$.  Equivalently, $H$ must be the union of conjugacy classes of $G$.

    But by inspection, the union of conjugacy classes of $G$ (including $\{1\}$) cannot have a number of elements that is a divisor of $60$ (other than $1$ or $60$).  So $H$ cannot exist. ~ $\blacksquare$
\end{proof}

The crux of the above argument, to reiterate, was the following useful observation:

\begin{theorem}{Normals are Conjugacy Classes}
    If $H \unlhd G$, then $H$ is the union of conjugacy classes of $G$, including the conjugacy class consisting of just the identity.
\end{theorem}

\begin{remark}
    It turns out that $G \cong A_5$.  One way to see this is to note:
    \begin{enumerate}
        \item The rotations $G$ of a regular icosahedron are isomorphic to the rotations $G'$ of a regular dodecahedron.
        \item There are $5$ cubes inscribed in a regular dodecahedron.
    \end{enumerate}
    Thus, $G$ is isomorphic to permutations of those $5$ cubes, and it turns out that these permutations must be even. \vspace{-0.2cm}
    \begin{center}
        \frame{\includegraphics[width=0.2\textwidth]{lecture-14/dodecahedron-via-cube}}
    \end{center}
    Note that point \#2 also gives a good strategy for drawing a dodecahedron\dots
\end{remark}

\subsection{Composition Series}

\textbf{Definition (Composition Series).} Let $G$ be a group.  Then a \textbf{composition series} of $G$ is a chain:
\[ G = G_0 \supseteq G_1 \supseteq G_2 \supseteq \dots \supseteq G_r = \{1\} \]
satisfying the property that for all $i$, we have $G_i \unlhd G_{i - 1}$ and $G_{i - 1} / G_i$ is simple.

\begin{remark}
    Confusingly, smaller groups get larger indices.  Also, simple groups are not trivial, so $G_{i - 1} \neq G_i$.
\end{remark}

\textbf{Example.} The group $C_6 = \langle x \rangle$ has two composition series:
\[ \langle x \rangle \supseteq \langle x^3 \rangle \supseteq \{1\} ~~ \text{ and } ~~ \langle x \rangle \supseteq \langle x^2 \rangle \supseteq \{1\}  \]
Meanwhile, the infinite group $(\mathbb{Z}, +)$ has no composition series: nontrivial subgroups of $\mathbb{Z}$ are isomorphic to $\mathbb{Z}$.

\begin{theorem}{Existence of Composition Series}
    Every finite group $G$ has a composition series.
\end{theorem}

\begin{proof}
    It's good enough to build the following \emph{chisel}.
\begin{boxed-text}[32em]
    \textbf{Chisel.} Given groups $G$ and $H$ with $H \unlhd G$ and $G/H$ not simple, \\
    
    the chisel constructs a group $T$ such that $H \unlhd T \unlhd G$, with $H \neq T \neq G$.
\end{boxed-text}
    To build a composition series given a chisel, we can just start with $\{1\} \unlhd G$ and use the chisel repeatedly.

    How does the chisel work?  Well, since $G / H$ is not simple, there must exist some normal subgroup $N \unlhd G / H$ besides $\{1\}$ and $G / H$.  So the chisel decides to choose $T := \pi^{-1}(N)$, where $\pi: G \to G / H$ is the quotient map.
    
    We now check the chisel's work by showing this choice of $T$ is valid:
    \begin{itemize}
        \item We know $H$ is a subgroup of $T$ and $T$ is a subgroup of $G$ both by the Correspondence Theorem.
        \item The fact $H \unlhd T$ follows from directly combining $H \unlhd G$ and $T \subseteq G$ together.  Indeed, \vspace{-0.1cm}
        \[ \text{``For all } g \in G, \text{ we have } gHg^{-1} = H\text{.''} ~~~~~  \text{ implies } ~~~~~ \text{``For all } t \in T, \text{ we have } tHt^{-1} = H\text{.''} \] \vspace{-0.85cm}
        \item The fact $T \unlhd G$ follows from the property $\pi(T) = N \unlhd G / H$.  Indeed, for any $g \in G$,
        \[ \pi(gTg^{-1}) = \pi(g) \pi(T) \pi(g^{-1}) = \pi(g) N \pi(g)^{-1} = N \text{ since } N \unlhd G / H. \]
        Thus, $gTg^{-1} \subseteq \pi^{-1}(N) = T$ for all $g$.  And $gTg^{-1} \subseteq T$ for all $g \in G$ is good enough to show $T \unlhd G$.
    \end{itemize}
    So the chisel works correctly, and we're done. ~ $\blacksquare$
\end{proof}

\subsection{The Jordan-Holder Theorem}

It turns out composition series are unique, up to order of ``prime factors''!

\begin{theorem}{Jordan-Holder}
    Suppose $G$ has two composition series:
    \[ G = A_0 \supseteq A_1 \supseteq \dots \supseteq A_r = \{1\} ~~ \text{ and } ~~ G = B_0 \supseteq B_1 \supseteq \dots \supseteq B_s = \{1\} \]
    Then the simple quotients $\{A_i / A_{i + 1}:  i = 0, 1, \dots, r - 1\}$ map bijectively with $\{B_i / B_{i + 1} :  i = 0, 1, \dots, s - 1\}$.
\end{theorem}

\textbf{Example.} The two composition series of $C_6$ have the same sets of simple quotients: $\{C_3, C_2\}$ and $\{C_2, C_3\}$.
\begin{align*} & \langle x \rangle \supseteq \langle x^3 \rangle \supseteq \{1\} ~~ \text{ has quotients } ~~ \langle x \rangle / \langle x^3 \rangle \cong C_3 \text{ and } \langle x^3 \rangle / \{1 \} \cong C_2. \\
& \langle x \rangle \supseteq \langle x^2 \rangle \supseteq \{1 \} ~~ \text{ has quotients } ~~ \langle x \rangle / \langle x^2 \rangle \cong C_2 \text{ and } \langle x^2 \rangle / \{1 \} \cong C_3. \end{align*}
We'll spend the rest of lecture proving the Jordan-Holder Theorem.  But first, some precautions:
\begin{itemize}
    \item Two non-isomorphic groups can have the same sets of simple quotients.  For example,
    \begin{align*} & S_3 \supseteq A_3 \supseteq \{1\} ~~ \text{ has quotients } ~~ S_3 / A_3 \cong C_2 \text{ and } A_3 / \{1\} \cong C_3. \\
    & C_6 \supseteq C_3 \supseteq \{1\} ~~ \text{ has quotients } ~~ C_6 / C_3 \cong C_2 \text{ and } C_3 / \{1\} \cong C_3. \end{align*}
    But unfortunately, $C_6 \not \cong S_3$.
    \item The condition $G_i \unlhd G_{i - 1}$ does \emph{not} imply, transitively, that $G_m \unlhd G_n$ for all $m \geq n$.  For example, taking
    \[ G_1 = S_4 \hspace{0.8cm} G_2 = \{1, ( 1\ 2)(3 \ 4), (1 \ 3)(2 \ 4), (1 \ 4)(2 \ 3)\} \hspace{0.8cm} G_3 = \{1, (1 \ 2)(3 \ 4)\}, \]
    we see that $G_2 \unlhd G_1$ and $G_3 \unlhd G_2$ both hold, but $G_3 \unlhd G_1$ does not hold.
\end{itemize}

\subsection{The Second Isomorphism Theorem}

We'll need two lemmas to prove the Jordan-Holder Theorem.  This one's the first.

\begin{theorem}{The Second Isomorphism Theorem}
    Suppose $H, K \subseteq G$ are subgroups with $K \unlhd G$.  Then:
    \begin{itemize}
        \item The set $HK = \{hk \mid h \in H, k \in K\}$ is a subgroup of $G$.
        \item The isomorphism $HK / K \cong H / (H \cap K)$ holds.  (Here, we are also implicitly claiming that $(H \cap K) \unlhd H$.)
    \end{itemize}
\end{theorem}

\begin{proof}
    To prove the first point, we simply check that $1 \in HK$ (obviously), that $(hk)^{-1} = h^{-1}(hk^{-1}h^{-1}) \in HK$, and:
    \[ hk h' k' = \underbrace{(hh')}_{\in H} \ \underbrace{(h'^{-1}k h')}_{\in K} \ \underbrace{(k')}_{\in K} \in HK. \]
    For the second point, we'll use the quotient map $\pi: G  \to G / K$.  Then:
    \begin{itemize}
        \item The First Isomorphism Theorem says $HK / K \cong \pi(HK)$.
        \item Since $K = \ker \pi$, we have $\pi(HK) = \pi(H)$.
        \item Well, $\pi(H)$ is itself the image of $\pi_{|H}: H \to G / K$.  But $\ker(\pi_{|H}) = H \cap K$, so $\pi(H) \cong H / (H \cap K)$.
    \end{itemize}
    Stringing together all three parts of the above yields $HK / K \cong H / (H \cap K)$, as desired. ~ $\blacksquare$
\vspace{0.1cm}
\begin{center}
    \frame{\includegraphics[width=0.4\textwidth]{lecture-14/second-isomorphism}}
\end{center}
\end{proof}

Some comments on the Second Isomorphism Theorem:
\begin{itemize}
    \item[\textbf{(\#1)}] Sometimes the Second Isomorphism Theorem is called the Third Isomorphism Theorem.
    \item[\textbf{(\#2)}] For an alternative proof of the first point, rephrase the first point more naturally like so:
    \begin{center}
        If $K \unlhd G$, then the group generated by $H$ and $K$ is $HK = \{hk \mid h \in H, k \in K\}$.
    \end{center}
    Usually, the group generated by $H$ and $K$ looks like $\{h_1k_1h_2k_2\dots \mid h_i \in H, k_i \in K\}$.  However, when $K \unlhd G$, we have the property that for any $h \in H$ and $k \in K$, there is some $k' \in K$ such that $kh = hk'$.
    
    So we may repeatedly use this property to shift all the $k_i$'s in the product $h_1k_1h_2k_2\dots$ to the right.
    \item[\textbf{(\#3)}] Nothing in $G$ matters besides elements in $HK$.  Thus, instead of mandating $K \unlhd G$, we could have gotten away with only mandating that $hKh^{-1} = K$ for all $h \in H$.

    In other words, the clause ``$K \unlhd G$'' in the theorem statement could have been replaced with ``$H \subseteq N(K)$''.  Here, $N(K)$ denotes the \textbf{normalizer} of $K$: the set $N(K) := \{g \in G: gKg^{-1} = K\}$.
\end{itemize}

\subsection{The Diamond Lemma}

Now for the second of the two lemmas.

\begin{theorem}{Diamond Lemma}
    If $H$ and $K$ are distinct normal subgroups of $G$ with $G / H$ and $G / K$ both simple, then \vspace{-0.05cm}
    \[ G / H \cong K / (H \cap K) ~ \text{ and } ~ G / K \cong H / (H \cap K). \] \vspace{-0.8cm}
\end{theorem}

\begin{proof}
    By symmetry, it suffices only to prove $G / K \cong H / (H \cap K)$.  Well, recall that the Second Isomorphism Theorem says $H / (H \cap K) \cong HK / K$.  So it suffices to prove $G / K \cong HK / K$.

    Since $H$ and $K$ are normal subgroups of $G$, so is $HK$.  Then $HK \unlhd G$ implies $(HK / K) \unlhd (G / K)$ as well.  But $G / K$ was assumed to be simple!  Thus, either $HK / K \cong \{1\}$ or $HK / K \cong G / K$.
    
    The latter is what we want, so assume the former: $HK / K \cong \{1\}$, or equivalently, $H \subseteq K$.  This lets us talk about $K / H$; in particular, $(K / H) \unlhd (G / H)$ holds.  But we now have another problem: $G / H$ was assumed to be simple!
    
    Thus, either $K / H = \{1\}$ or $K / H = G / H$.
    \begin{itemize}
        \item The former implies $K = H$, contradicting the assumption $K \neq H$ in the theorem statement.
        \item The latter implies $K = G$, contradicting the assumption that $G / K$ was simple.  (Trivial groups aren't simple!)
    \end{itemize}
    So we're done. ~ $\blacksquare$
\end{proof}

There's not much to say about the Diamond Lemma.  The following image reveals the meaning behind its name:
\begin{center}
    \frame{\includegraphics[width=0.4\textwidth]{lecture-14/diamond-lemma}}
\end{center}

\subsection{Proof of The Jordan-Holder Theorem}

We're now ready to prove the Jordan-Holder Theorem.  Recall that we have a group $G$ with two composition series:
    \[ G = A_0 \supseteq A_1 \supseteq \dots \supseteq A_r = \{1\} ~~ \text{ and } ~~ G = B_0 \supseteq B_1 \supseteq \dots \supseteq B_s = \{1\} \]
We'd like to show these two series are \emph{equivalent}; that is, their simple quotients $\{A_i / A_{i + 1}\}$ and $\{B_i / B_{i + 1}\}$ map bijectively with each other.  Our approach is induction on $r$, with inductive hypothesis:
\begin{center}
    \textbf{Inductive Hypothesis.} If a group $G'$ has \emph{any} composition series of length $r - 1$, \\ then all composition series of $G'$ are equivalent.
\end{center}
(Here we WLOG $r \geq s$.) The entire rest of the proof is captured in the following chain of equivalences:

\begin{center}
    \frame{\includegraphics[width=0.7\textwidth]{lecture-14/jordan-holder-proof}}
\end{center}

There are two steps to the proof: step \circled{1} and step \circled{2}.

\begin{enumerate}
    \item[\circled{1}] We may assume $A_1 \neq B_1$; indeed, if $A_1 = B_1$, then the inductive hypothesis would tell us $A_1 \supseteq \dots \supseteq A_r$ and $B_1 \supseteq \dots \supseteq B_s$ were equivalent, and we'd be done.
    
    Now consider $K_1 = A_1 \cap B_1$, along with \emph{any} composition series $K_1 \supseteq \dots \supseteq K_t$ of $K_1$.
    \begin{quote}
        \textbf{Claim.} The chain $A_1 \supseteq K_1 \supseteq \dots \supseteq K_t$ is a valid composition series of $A_1$.

        \emph{Proof:} It suffices to show $A_1 / K_1$ is simple.  This follows from the Diamond Lemma, which says $A_1 / K_1 \cong G / B_1$.  And $G / B_1$ is simple by definition. ~ $\square$
    \end{quote}
    Thus, we have constructed \emph{two} composition series of $A_1$ (green in the original diagram):
    \[ A_1 \supseteq A_2 \supseteq \dots \supseteq A_r ~~ \text{ and } ~~ A_1 \supseteq K_1 \supseteq \dots \supseteq K_t. \]
    But the former of these two composition series has length $r - 1$.  So by the inductive hypothesis, these two composition series are equivalent, as desired.
    \item[\circled{2}] This follows \emph{immediately} from the Diamond Lemma, which says $G / A_1 \cong B_1 / K_1$, and $G / B_1 \cong A_1 / K_1$.
\end{enumerate}
And that's the proof! ~ $\blacksquare$

\end{document}